\section{La préambule}\label{la-preambule}

\subsection{La mission}\label{la-mission}

Un des principaux mandats de la Fédération canadienne étudiante de génie (FCEG) est de défendre les intérêts de plus de 80 000 étudiants en ingénierie répartis dans 45 établissements d'enseignement supérieur à travers le Canada. Au fil des ans, la FCEG a établi des partenariats avec des associations telles que le Bureau canadien d'accréditation des programmes de génie (BCAPG) et les Doyennes et doyens d'ingénierie Canada (DDIC), qui ont permis à la FCEG de défendre les intérêts des étudiants en génie, d'influencer la prise de décision concernant le programme d'études en génie et de façonner l'avenir de la profession d'ingénieur.

\subsection{Le portefeuille d'activités de plaidoyer de la FCEG, en bref}\label{le-portefeuille-dactivites-de-plaidoyer-de-la-fceg-en-bref}

Les activités de plaidoyer de la FCEG reposent sur deux (2) activités principales : les prises de position nationales et les enquêtes nationales. Les prises de positions sont destinées à être mises à jour régulièrement pour refléter les avancées et les reculs dans la défense des intérêts des étudiants en ingénierie et sont largement basées sur les données collectées dans le cadre des enquêtes nationales auprès des étudiants, qui sont censées être publiées une fois par an.

La FCEG a publié trois (3) enquêtes nationales à ce jour : l'édition de 2017, l'édition de 2020 portant sur des thèmes spécifiques à la COVID-19, et l'édition de 2023 sur le développement et le soutien des étudiants. La FCEG a adopté ses premières positions lors d'une assemblée générale à l'occasion de la Conférence canadienne sur le leadership en ingénierie (anciennement connu sous le nom de Congrès) en 2018, et elles ont depuis été mises à jour en 2024.

\subsection{Le cahier des positions}\label{le-cahier-des-positions}

Le cahier des positions récapitule les prises de positions officielles de la FCEG approuvées par ses membres, afin de mettre en avant les efforts des activités de plaidoyer dans divers aspects de l'expérience des étudiants en ingénierie. Le cahier des positions et les prises de position individuelles sont mis à la disposition des membres dans le cadre de leurs propres activités de plaidoyer. Ces documents s'inspirent des prises de position d'Ingénieurs Canada et constituent une ressource précieuse pour les membres, qui peuvent s'en inspirer pour leurs initiatives de plaidoyer institutionnelles.

\newpage
